两层无界横截面平板A和B各厚 \(L\)，互相完美热接触。一维定常导热，A层热导率 \(\kappa_A(T)=aT\)，B层 \(\kappa_B(T)=b\)，其中 \(a,b>0\)，无内热源，忽略界面接触热阻。两端热库分别为 \(T_h>T_c>0\)，固定两层的先后次序A在左B在右。
(1) 将Fourier定律分段积分，写出左热右冷时的界面温度方程与热流。
(2) 当右热左冷时重新求界温及热流，证明一般两个方向热流大小不相等，并明确必须保持同一材料排列。
(3) 以无量纲 \(L=a=b=1,T_h=2,T_c=1\) 为实例求两个方向精确界温、热流与整流比，核对其大小关系。
(4) 由各层局部熵产生 \(\kappa(T)(T_x)^2/T^2\) 积分得到总熵生产，验证热库的总熵账本并讨论 \(a\to0\) 和零温差极限。
